\documentclass[11pt]{article}
\usepackage[a4paper,top=2.0cm,bottom=2.0cm,left=2.2cm,right=2.2cm]{geometry}
\usepackage[T1]{fontenc}
\usepackage{lmodern}
\usepackage[english]{babel}
\usepackage{microtype}
\usepackage{amsmath,amssymb}
\usepackage[hidelinks]{hyperref}
\usepackage{enumitem}
\usepackage{parskip}
\setlength{\parskip}{0.4em}

\newenvironment{referee}{\begin{quote}\itshape\small}{\end{quote}}
\newcommand{\reply}{\par\noindent\textbf{Reply.}\ }
\newcommand{\dd}{\mathrm{d}}
% ---------------------------------------------------------------
% GENERATED FILE -- do not edit by hand.
% Cross-reference numbers read out of paper2.aux by
% provenance/make_crossrefs.py, so that the response cannot cite a
% number the manuscript does not have.
% ---------------------------------------------------------------
\newcommand{\msSecDomain}{2.1}
\newcommand{\msSecBvp}{3.4}
\newcommand{\msSecClassification}{3.2}
\newcommand{\msSecInversion}{3.3}
\newcommand{\msSecThakurta}{2}
\newcommand{\msSecSubmersion}{2.2}
\newcommand{\msAppValidation}{Appendix A}
\newcommand{\msLemCompact}{2}
\newcommand{\msProtStatus}{1}
\newcommand{\msProtEndpoint}{2}
\newcommand{\msProtNames}{3}
\newcommand{\msProtEvidence}{4}
\newcommand{\msPropClassification}{4}
\newcommand{\msPropExteriorSep}{5}

% ---------------------------------------------------------------
% GENERATED FILE -- do not edit by hand.
% Produced by paper2/provenance/make_provenance.py; see that file for
% why these numbers are generated rather than transcribed
% (CQG-116884, referee major comments 7 and 10).
% generated 2026-09-23T23:15:40+00:00
% content digest (reproduces on any machine; timings excluded): 42ed5ab1779c58132ab9a4349e3f4bf1c28b40907b883efb7d1ed83af0b4ba25
\newcommand{\ProvenanceDigest}{\texttt{42ed5ab1779c5813}}
\newcommand{\ProvenanceRelease}{v1.7.1}
\newcommand{\ProvenanceDOI}{10.5281/zenodo.22926710}
\newcommand{\ProvenanceConceptDOI}{10.5281/zenodo.21707377}
\newcommand{\ManifestArtefacts}{31}
% ---------------------------------------------------------------
\newcommand{\provEpsWindow}{$\varepsilon\in\{0.001,\,0.002,\,0.004,\,0.008,\,0.016\}$}
\newcommand{\provFitInterval}{$r/M\in[5.43,10.41]$ across the three configurations}
\newcommand{\provLocusSep}{1.93\,M}
\newcommand{\provBindingLocus}{moving turning point}
\newcommand{\provMinHpr}{$0.352$}
\newcommand{\provMinVbarSq}{$0.483$}
\newcommand{\provExactSlope}{$2.12\pm0.03$}
\newcommand{\provLeadSlope}{$1.01$--$1.02$}
\newcommand{\provExactRunSlope}{$2.15$}
\newcommand{\provExactRunResidual}{$1.90\times10^{-6}$}

% per-configuration rows of the validation run
\newcommand{\provValidationRows}{%
  $t$, $a{=}0.9,\hat E{=}1.4,J{=}6$ & $1.01\pm0.00$ & $2.12\pm0.03$ \\
  $\tau$, $a{=}0.9,\hat E{=}1.4,J{=}2.5$ & $1.02\pm0.01$ & $2.12\pm0.03$ \\
  $t$, $a{=}0.5,\hat E{=}1.3,J{=}5$ & $1.02\pm0.01$ & $2.12\pm0.03$ \\
}

% fitted exponent on shrinking epsilon sub-windows
\newcommand{\provSubwindowDrift}{$2.118 \to 2.068 \to 2.043$}
\newcommand{\provSubwindowNarrowest}{$2.045$}
\newcommand{\provSubwindowEpsMax}{0.004}

% SHA-256 (first 12 hex) of every script behind the table
\newcommand{\provHashFigPhiValidationCorrected}{\texttt{476e662efce9}}
\newcommand{\provHashAdiabaticFirstOrderExact}{\texttt{3cd72fd798e3}}
\newcommand{\provHashCuspideErgosfera}{\texttt{55fc540ecf17}}
\newcommand{\provHashAdiabaticConvergenceRaw}{\texttt{758b59f23295}}

\newcommand{\provHashLine}{%
  SHA-256 (first 12 hex): \texttt{fig\_phi\_validation\_corrected.py}\,$=$\,\texttt{476e662efce9};
  \texttt{adiabatic\_first\_order\_exact.py}\,$=$\,\texttt{3cd72fd798e3};
  \texttt{cuspide\_ergosfera.py}\,$=$\,\texttt{55fc540ecf17};
  \texttt{adiabatic\_convergence\_raw.npz}\,$=$\,\texttt{758b59f23295}.}

\newcommand{\R}[1]{\par\medskip\noindent\textbf{#1}\par\nopagebreak}

\title{\vspace{-1.6cm}Response to the referees\\[2pt]
\large Manuscript CQG-116884}
\author{Iman Rosignoli}
\date{\today}

\begin{document}
\maketitle
\thispagestyle{empty}

\noindent I thank both referees. The first report was unusually specific: it found a
defect I had not seen, which turned out to be more serious than the report could
establish from outside, and correcting it produced a stronger result than the one it
replaced. The second sent me to a body of geometry and to an experimental question I
had not engaged with, and the paper now contains work that exists because of it.

The revision follows Referee~1's reconstruction list in the order given.

\begin{enumerate}[leftmargin=1.4em,itemsep=1pt]
\item The exterior/contact/continuation distinction is enforced globally. New
section~\msSecDomain\ proves as Lemma~\msLemCompact\ the compactness criterion the
referee sketched, from which the compact-control problem \emph{ends} at $r=2M$ for
$W=\partial_\eta$; Protocol~\msProtStatus\ fixes the labelling (E)/(C)/(A) and is
applied throughout, figures included.
\item The extremal-versus-minimiser statements are corrected in both papers, and the
claim that numerical perturbation establishes minimality is gone.
\item The $J=-J_c$ contradiction is resolved by rederivation. \emph{Both} printed
statements were wrong; the corrected result is a four-case classification
(Proposition~\msPropClassification), and all affected material is regenerated.
\item The archive is rebuilt: the missing module restored, slopes and checksums
generated rather than transcribed, and a manifest of \ManifestArtefacts\ entries
mapping artefacts to commands and checksums.
\end{enumerate}

\noindent The exterior retrograde separatrix is now the flagship exact result, as
recommended, and the title no longer says ``trichotomy''.

\bigskip\hrule
\section*{Referee 1}

\R{Major 1. The control domain ends at the conformal stationary limit.}
\reply Accepted in full; this is the structural change of the revision.
Lemma~\msLemCompact\ states that the rail slice is nonempty and compact iff $W$ is
timelike and $\hat E\ge|W|$, degenerating to a point at $\hat E=|W|$ and empty or
noncompact otherwise. Since $g(\partial_\eta,\partial_\eta)=-A^2(1-2M/r)$, the two
conditions read $r>2M$ and $\bar v^2=1-A^2f/\hat E^2>0$. The text says explicitly
that a change of chart does not repair this---Doran coordinates regularise the
formulae at $\Delta=0$ but leave $g(W,W)$ unchanged---and that a Doran- or
ZAMO-anchored selector would be a different problem, not treated here.

A precedent worth recording came out of the second report's reading list.
Giannoni, Piccione and Verderesi pose the stationary problem with $-\langle\dot z,Y\rangle\equiv k>0$ and
restrict at once to $|Y|<k$, which is Lemma~\msLemCompact\ exactly. The restriction
has stood since 1997; what is new here is its behaviour once $W$ fails to be
Killing, when the boundary moves with $A(\eta)$.

\R{Major 2. Pontryagin extremals are not automatically local minima.}
\reply You are right, and the error was worse than reported: the sentence was in
\emph{both} versions of Paper~I. Both now state the hierarchy as you write it,
minimiser $\Rightarrow$ PMP extremal and PMP extremal $\nRightarrow$ local
minimiser, with the note that local minimality needs a second-variation or
conjugate-point argument.

The revision also supplies sufficient conditions where they hold. On the frozen
arrival-time branch the optical geometry is Riemannian with Gauss curvature negative
throughout the exterior exactly when $\hat E^2\ge3/2$; there are then no equatorial
conjugate points, and by the index theorem in the form of Giannoni, Masiello and
Piccione (\emph{J. Geom. Phys.} \textbf{35} (2000) 1--34) every exterior extremal is
a local minimiser. The surface is complete, so Cartan--Hadamard on the universal
cover gives one global minimiser per winding class, and the optical distance
saturates the eikonal equation, supplying the value function your list mentions.

These statements are equatorial by nature, not by present limitation, and the
manuscript says so. The tangential curvature $K_{\rm tan}=1/B>0$ at the circular
orbit, so Cartan--Hadamard cannot apply in three dimensions; and the winding classes
themselves do not survive the lift, the equatorial surface being
$(2M,\infty)\times S^1$ while the full one is simply connected. The out-of-plane
Jacobi field is $\sqrt{B}\,|\sin\varphi|$, which vanishes at swept azimuth $\pi$;
for a fixed non-radial arc, $K_t<0$ together with $\Phi<\pi$ is therefore
\emph{sufficient} for non-conjugacy in every direction, and nothing stronger
follows. In particular I do not import the $\Phi<\pi$ bound from Paper~I, which is
proved there for the proper-time family, whereas these statements are on the
arrival branch. Nothing analogous is claimed on the proper-time branch, where
$r^3K_\tau\to+M>0$ and the surface is incomplete at the stationary limit.

The passage you will have had in mind, on perturbing the computed brachistochrones,
is gone; the appendix now says such perturbations are a consistency check and cannot
certify a minimum in an infinite-dimensional class.

\R{Major 3. Endpoint protocol and costate relabelling.}
\reply Accepted. Protocol~\msProtEndpoint\ fixes the default---fixed launch, fixed
spatial target, free arrival clock, branch-dependent shooting---and comparisons
between branches are made at $J_t=A\,J_\eta$, so ``same $J$'' and ``same endpoints''
are different comparisons. Section~\msSecBvp, the one place where the default does
not apply, now names its two protocols, records that at a common numerical $J$ they
compare different orbits, and adds that an ordering between two selected extremals
does not show either is the global minimiser. The algebraic regularity of the
printed Hamiltonians at $f=0$ is separated from the control-domain question.

\R{Major 4. The retrograde marginal orbit: cusp or corner.}
\reply Accepted, and this was the most useful comment in the report. I rederived
the case as instructed, and \emph{both} printed statements were wrong.

Two exact identities organise it: $p_r^2=\mathcal S/(D_E^2\Delta^2)$ with the
prefactor regular and nonzero at $r_e$, and at the marginal value the exact
factorisation $r\Delta-J_c^2D_E=(r-2M)(r^2+J_c^2)$, giving
$\mathcal S=r^4f^2w(r^2+J_c^2)$. So $r_e$ is a double root for \emph{both} signs and
reality is not what selects the prograde value; the previous counter-rotation
argument is withdrawn, and it was wrong in detail too, since near $r_e$ the
retrograde marginal has $\dot\varphi>0$ and is locally co-rotating.

What breaks the $\pm J$ symmetry is invisible to the radical and visible to the
Hamilton equations: $\tilde p_\varphi(r_e)=J-J_c$ is linear in $J$, and
$\mathcal T(r_e)=-a(J-J_c)/\bar P_e$ vanishes only at $J=+J_c$. The classification
is fourfold: smooth periapsis for $|J|>J_c$; semicubical cusp for $|J|<J_c$;
at $J=-J_c$ a double root with $r_e$ reached \emph{asymptotically only}, finite
tangent, exponent~1; and a regular crossing at $J=+J_c$. So the main text's ``cusp''
was wrong and the appendix's ``corner'' half right---the tangent is finite, but
there is no reflection. Your caution against merging the generic and marginal cases
is justified quantitatively: the cusp amplitude $K=(EJ/a^2)\sqrt{2M/(a^2-E^2J^2)}$
\emph{diverges} as $|J|\to J_c^-$.

The one-sided expansions, signs, branch orientation and convergence tests are in
\texttt{verify\_cusp\_corner.wls} (21 checks) and
\texttt{verify\_marginal\_hamilton.wls} (19 checks), all passing. At $M=1,a=0.9,E=1.2$
the retrograde rate is $C=-0.267000117$ against the closed form to $10^{-6}$; the
elapsed parameter gains $34.4956$ per four decades against $\ln(10^4)/|C|=34.4955$.
All affected material was regenerated, not recaptioned: the atlas had drawn a
reflection at $-J_c$, and Table~A1 is now keyed on root multiplicity and attainment.

\R{Major 5. The exterior separatrix as flagship result.}
\reply Accepted and adopted as the organising result.
Proposition~\msPropExteriorSep\ records $r_d=3.513905124011657\,M$ and
$J_c^-=-8.053516003877019$, reproducing your values to sixteen digits, with the
printed closed form agreeing to $10^{-12}$. At $r_d$ both control-domain
inequalities hold strictly, $g(W,W)=-0.430833$ and $\bar v^2=0.700811$, and they
hold on the whole approach arc, so it is type~(E) rather than a continuation. A spin
scan gives $r_d/M=2.857675,\,3.204756,\,3.513905,\,3.579494$ at
$a=0.1,\,0.5,\,0.9,\,0.99$, so the exterior character is structural. The prograde
degeneration at $r_d=1.512292M$ lies outside $r_+=1.435890M$ but inside $2M$, and is
recorded as type~(A), not a physical control separatrix. This also answers M4.

\R{Major 6. Tiered proof status of the fixed-endpoint result.}
\reply Accepted. The conclusions keep five levels apart: an analytic integrand
inequality; a theorem for $r_0\le R_*(E,a)$; an asymptotic result as
$r_0\to\infty$ in the static case; a computer-assisted proposition on a stated grid;
and a conjecture in general. Protocol~\msProtEvidence\ gives them different
evidential status. The text also says that an ordering between two selected
extremals does not prove either is the global minimiser. On executability you were
right that the CAP runner imported a module that was absent; see major~10.

\R{Major 7. Scope and provenance of the first-order result.}
\reply Accepted, and the problem was worse than the discrepancy you saw: the
numbers had been transcribed by hand from separate runs, and re-running gave
$2.12\pm0.03$ and $2.15$, so neither printed value was current. They are now
generated. One archived command,
\texttt{python3 paper2/provenance/make\_provenance.py}, runs each validation script
as a subprocess, reads the $\varepsilon$-window from the source, hashes every script
and writes the machine-readable record, the \LaTeX{} macros the manuscript
\texttt{\textbackslash input}s, and the manifest.

The theorem is now local on a compact regular subarc, with turning points, the
freezing surface, the stationary limit, chart singularities and moving-separatrix
crossings excluded, the last identified as a singular inner problem. On the numbers:
since the residual is $c_2\varepsilon^2(1+k\varepsilon+\cdots)$ the local exponent
is $2+O(\varepsilon)$, so a finite-window fit is biased upward and $\sigma$ measures
fit quality rather than that bias. Refitting on shrinking sub-windows gives
$2.118\to2.068\to2.043$, $2.123\to2.072\to2.045$ and $2.122\to2.070\to2.044$ in the
three configurations---every one decreasing toward $2$. The window is
\provEpsWindow\ and the fitted interval \provFitInterval, $6M$ from the nearest
excluded locus.

\R{Major 8. Higher-genus terminology.}
\reply Accepted. A convention now states that \emph{length-two iterated Abelian
integral on the genus-two hyperelliptic curve} is the only classification claimed,
and that ``genus-two dilogarithm'' is descriptive shorthand for it, not a claim that
a function class with specified kernels, relations, monodromy or single-valued
completion has been established. The three theorem-level bare uses are converted.
``Elliptic dilogarithm'' is retained, since at genus one the object is established.

\R{Major 9. Physical language.}
\reply Accepted. Protocol~\msProtNames\ fixes the dictionary and is applied
throughout: $\hat E$ is the actively maintained rail charge, not a conserved energy;
$J=p_\varphi$ is the axial control costate, not mechanical angular momentum; the
rail minimises a selected clock, not thrust or fuel; $r_e=2M$ is the conformal
stationary limit; $r_+$ is the seed Kerr null surface, with the note that for
time-dependent $A(\eta)$ neither horizon has been computed; and inside $2M$ we write
``continued orbit'', ``root pattern'' or ``analytic branch''.

\R{Major 10. Reproducibility package.}
\reply All three findings are correct, and I am glad they were checked.
\texttt{cap\_full.py} had been factored out during development and never archived;
it is restored verbatim from the prototype, with a self-test. On the sibling
directories the diagnosis turned out to be different: all thirteen generators were
tracked, but each inserted the repository root on \texttt{sys.path} while the shared
style module lives one directory down, so on a fresh clone every one failed at
import. Fixed in all thirteen---which also exposed a silent bug, a turning-radius
helper returning the grid point just inside the forbidden region, dropping one curve
from a panel with no error raised. The printed hashes were stale; the same two files
carried three different values, which is the signature of hand transcription, and
they are now generated by the command above.

Two scope statements belong with this. The content digest covers the
adiabatic-validation scripts, windows and fitted data listed in it; it is not a
whole-archive checksum, and each manifest row separately carries the checksum of its
own file. The manifest maps the principal computational claims and the figures and
tables depending on them, not every equation.

\R{Major 11. Dependence on Paper I.}
\reply Accepted. The incorrect PMP sentence is corrected there. Paper~I's existence
and normality theorem now defines the extended admissible class and clock state,
states the regularity assumptions, explains closure of the non-autonomous clock
equation, and identifies the endpoint protocol and excluded degeneracies. The Vaidya
qualification is carried into the abstract and conclusion: the compact-control
domain ends at $r=2m(v)$, the calculation establishes an exterior approach and
contact threshold rather than an interior optimal trajectory, and $m'(v)<0$ in the
ingoing metric is a formal continuation. Paper~II no longer imports optimality
claims from Paper~I; it cites the verification criterion as a criterion.

\R{Major 12. Relation to established work.}
\reply Accepted, and treated as positioning. A new introduction paragraph separates
established from new. Established: Perlick's fixed-energy optical metric; the
sub-Riemannian formulation of Giannoni, Piccione and Verderesi; the arrival-time
theory of Giannoni and Piccione; the Morse theory of Giannoni, Piccione and Tausk. I
have added the Giannoni--Masiello--Piccione programme---including the finiteness
results on \emph{conformally stationary} spacetimes, the closest published setting to
the frozen-$A$ regime used here---and Caponio, Corona, Giamb\`o and Piccione on
fixed-energy solutions of an indefinite Lagrangian with an affine Noether charge.

The distinction is stated there: in all of those the charge is \emph{conserved},
because the generator is a symmetry; here it is held by \emph{control} against a
background of which $W$ is in general not a symmetry, at the cost of proper
acceleration, so $\hat E$ is a design parameter. Three consequences have no
stationary counterpart: possible loss of compactness and hence a causal boundary;
the reduced Hamiltonian ceasing to be an interior first integral; and first-order
response as a genuine perturbation problem. I have also taken up your closing
framing, putting the emphasis on the controlled invariant and the causal boundary
rather than on another exact trajectory.

\R{Minor comments.}
\begin{description}[leftmargin=2.4em,style=nextline,itemsep=0pt]
\item[M1] Done: Protocol~\msProtEndpoint, stated once and used consistently.
\item[M2] Done: $J=p_\varphi$ identified as a control costate at first use.
\item[M3] Done: the compact indicatrix boundary and the convex body whose support
function is used are distinguished, with the reason no convexification is needed.
\item[M4] Done: see major~5.
\item[M5] Done. A mechanical audit of all captions, re-run against the final text,
found five quoting radii or momenta without the normalisation; all now carry it.
\item[M6] Done: every figure continuing a curve inside $r_e$ draws that portion in a
distinct line style and states, in the legend or panel title, that the continuation
is analytic and no optimality is claimed. Figures whose curves terminate at $r_e$
carry no such label, because nothing is continued there.
\item[M7] Done: Protocol~\msProtNames\ carries a table keeping $E$, $\hat E$,
$E_{\rm eff}$, $J$ and $J_{\rm eff}$ distinct, including $J_t=A\,J_\eta$.
\item[M8] Done: deviations declared absolute unless labelled otherwise, norms named
where the residual is a vector, every fitted slope quoted with its interval.
\item[M9] Done: Protocol~\msProtEvidence\ separates the evidence levels and claims
are tagged.
\item[M10] Done: the conclusion is shortened and lists only results whose domain and
proof status are established in the body.
\end{description}

\bigskip\hrule
\section*{Referee 2}

\R{Bibliography.}
\reply Accepted with thanks; the list was useful. I have added or made explicit
Giannoni, Piccione and Verderesi (1997); Giannoni and Piccione on arrival-time
brachistochrones; Giannoni, Piccione and Tausk, with the \texttt{math-ph/9905007}
identifier so that preprint and published version are visibly the same work; Haws
and Kiser (1995); and Ta\c{s} (\texttt{arXiv:2512.08776}). Your suggestion of
connecting the global analysis to the Morse index has shaped the revision: it is the
framework in which Referee~1's extremal-versus-minimiser request is properly
answered.

One bibliographic note, so the record is right: the article at \emph{J. Math. Phys.}
\textbf{38}(12), 6367--6381 (1997) given as item~[6] is by Giannoni, Piccione and
Verderesi; Giannoni, Masiello and Piccione are a separate collaboration, and I now
cite three of their papers in their own right, including the finiteness results I had
missed. Checking your list also caught a wrong page number in my own entry.

\R{Curvature-operator eigenvalues and weighted manifolds.}
\reply This comment sent me back to the geometry and it was worth it: there is a
precise structural relation between the two frameworks, which the manuscript had left
implicit, and a new section~\msSecSubmersion\ now sets it out. The relation is not
quite the one the comment proposes, and the difference is one of \emph{scope} rather
than of correctness---the theorem holds where its hypotheses hold, and the objects
here fall outside them.

Three connections are now in the paper because of this report. First, the optical
reduction is a submersion, and with Perlick's metric on the base it is horizontally
conformal with dilation $\Lambda^2=\hat E^2/[f(\hat E^2-f)]$, the Perlick factor
itself. Second, since $\mathcal L_Wg=(2A'/A)g$, the selector is conformal Killing
and $\dd\hat E/\dd\tau|_{\rm geodesic}=A'/A=\varepsilon$ independently of the
four-velocity: the rate at which an uncontrolled worldline loses the rail is the
adiabatic expansion rate, which makes the middle rung of the
Killing~$\to$~conformal-Killing~$\to$~Kodama hierarchy quantitative. Third, choosing
the weight $e^{-2f^2}=A^{-2}$ makes the weighted metric \emph{exactly} the Kerr seed,
with $\operatorname{div}_{\bar g}W=0$ so that your equation~(32) holds verbatim,
while $\operatorname{div}_gW=4A'/A$; your Theorem~T8 requires a parallel weight, and
ours is parallel precisely at frozen $A$. That framework therefore describes the
frozen level of this construction exactly, and the instinct behind the comment is
well founded.

In revising I corrected my own statement of the first point. What I had called the
``Fuglede--Ishihara dichotomy'' is Fuglede's \emph{definition} of semiconformality
(\emph{Ann. Inst. Fourier} \textbf{28} (1978) 107, \S5), not a theorem, and invoking it
would have been circular. Pursuing Ishihara's version instead gave a definite answer.
His reduction of harmonicity to minimal fibres holds at unit dilation; ours has
non-constant dilation, so the tension field carries a second term,
$\tau(\pi)=-\dd\pi(\mu+\nabla^H\ln\Lambda)$, and non-geodesic fibres alone leave open
that the two cancel. Computing both, with $\mu=(M/r^2)\partial_r$, gives
\[
   \mu+\nabla^H\ln\Lambda=\frac{M}{r^2}\,\frac{f}{\hat E^2-f}\,\partial_r\neq0
   \qquad(r>2M,\ \hat E>1),
\]
strictly positive throughout the exterior. So the projection is not a harmonic
morphism, and the obstruction is the very quantity that makes the rail cost
something. The claim is the non-vanishing of a computed tension field on the
horizontal distribution, which is spacelike; I do not invoke the global elliptic
theory, which is Riemannian. The identity is checked in two independent
computer-algebra systems.

On Myers, my first argument was too quick---decay at infinity does not exclude his
hypothesis, and he supplies the counterexample himself. Tested on
$\mathrm{Ric}(v,v)$ instead, the optical scalar curvature is negative at every
exterior point and every rail energy, whereas $\mathrm{Ric}(v,v)\ge e^2>0$ would
force $R_{\rm opt}\ge3e^2>0$. The hypothesis fails by sign, pointwise, and the
relevant instrument is his Lemma, which needs no completeness. Correcting a non
sequitur produced a better result.

Set beside the objects treated here, Theorem~T2 is speaking about something else, in
three concrete ways. Equation~(9) is a Hermite-type reduction of abelian differentials
on a hyperelliptic curve; a dichotomy about the differential of a map and a reduction
of abelian integrals are statements about different objects. A controlled-rail
brachistochrone is not a geodesic of any weighted metric but an extremal of a
time-optimal control problem with an actively enforced pointwise constraint, a compact
indicatrix and a support Hamiltonian---and there is a regime where the control set
loses compactness and the problem ceases to be posed, which has no counterpart for
geodesics. And Remark~R4 stresses that the photon follows geodesics while
Theorem~T25 constructs a \emph{Killing} field from the four-velocity; the rail is
neither, which is why the three-rung hierarchy is needed at all.

For these reasons the results above do not follow from Theorem~T2. I want to be exact
about how narrow that is: it is not a judgement on the theorem, nor on the framework,
which I have adopted and credited. Of the connections this report proposed, three are
now in the manuscript, together with the Myers correction; the fourth, the implication
from T2, is the one I have not been able to make. If it does exist I would be glad to
cite it. If the editor thinks the connection worth recording, I would add a neutral
remark that weighted-manifold and curvature-operator methods provide an alternative
setting for related variational problems, without a dependence claim.

\R{Experimental validation.}
\reply Accepted, and the revision answers it rather than declining it. No
gravitational validation is available and none is claimed, but the request was for
techniques, and those exist for the parts of this work that are not the Einstein
dynamics. \msAppValidation\ sets out three levels.

The controlled rail is realisable on a bench: a holonomic planar robot or $XY$ stage,
tracked from above, commanded so that its \emph{measured} velocity obeys
$\dot x=c+Re(\theta)$. Applying 32--64 control directions on a grid of states and
fitting the velocity locus measures the indicatrix itself against the $c$ and $R$
predicted here, with no refitting of trajectories; a programmed $A(\chi)$ then tests
the first-order response and the $O(\varepsilon^2)$ scaling, and the bench can be
driven across $J=-J_c$ for the four-case behaviour. This validates the reduced system
and its control structure, not the field equations, and I do not present it
otherwise. A design study is in the archive under \texttt{Experimental/}, runnable as
one command, with the indicatrix identity closing to $4.4\times10^{-16}$ and the
frozen-flow Hamiltonian conserved to $\sim10^{-15}$. It is a design study, not data,
and is labelled so.

Analogue platforms supply the geometry. Rotational superradiance was measured on a
draining vortex---amplification $14\%\pm8\%$ at 3.70~Hz, $R_{m=1}=1.09\pm0.03$,
$R_{m=2}=1.14\pm0.08$ (Torres \emph{et al.}, \emph{Nature Phys.} \textbf{13} (2017)
833)---and the same structure has been realised in superfluid $^4$He
(\v{S}van\v{c}ara \emph{et al.}, \emph{Nature} \textbf{628} (2024) 66). Their
effective metric is a threading form with wind $\boldsymbol v$, hence Randers, and its
ergosurface sits at $|\boldsymbol v|=c$, that is at $g(W,W)=0$, where the selector
ceases to be timelike---the \emph{first} hypothesis of Lemma~\msLemCompact, and the
same clause that ends our control domain at $r=2M$. It is not the energy-dependent
freezing condition $\hat E=|W|$, at which $W$ is still timelike. The time-dependent
half has its own realisations, in expanding condensates (Eckel \emph{et al.},
\emph{Phys. Rev. X} \textbf{8} (2018) 021021) and in the programmable scale factor of
Viermann \emph{et al.} (\emph{Nature} \textbf{611} (2022) 260). The two features of
$g=A(\eta)^2g_{\rm Kerr}$ each have a laboratory realisation and no platform I know
of carries both; combining them is the experimental question this geometry poses.

Two limitations belong with that. These platforms carry \emph{waves}, so they probe
the geometry and its boundary rather than an actively controlled massive worldline;
and a purely conformal rescaling is invisible to ideal null rays, so an analogue sees
$A(\eta)$ through mode redshift and particle production rather than ray trajectories.
It is because the rail is timelike at fixed $\hat E$ that the conformal factor acts
on it at all---the same reason the problem does not reduce to Fermat's principle.

Since no measurement is on offer, the revision states in one place what an
independent recomputation must return: the separatrix at
$r_d/M=3.513905124011657$ with $J_c^-=-8.053516003877019$, the threshold
$\hat E^2=3/2$, negativity of the optical scalar curvature, the twist $-a/2M$, the
fibre acceleration, the minimum thrust, and the adiabatic exponents. Each is quoted
with enough digits that a sign error cannot hide, and each localises a different part
of the construction. That is the form of falsifiability this problem admits, and I
would rather commit to it than claim an observational contact the work does not have.

\bigskip\hrule
\section*{A note of thanks}

Referee~1's report caused four errors to be found---the cusp/corner contradiction,
the PMP statement, the absent module and the drifting hashes---and pursuing them
turned up three more the report could not have seen: a broken import path that
prevented every figure script from running on a fresh clone, a silent turning-radius
bug that dropped a curve from a published panel, and a wrong page number.

Referee~2's report is responsible for material that is new rather than corrected: the
reading list, which caught that page number and the finiteness results I had missed;
the insistence on experimental validation, which produced the three-level section;
and the geometric framework, which produced three connections now in the paper
together with the Myers correction. Where I have not been able to follow that report
I have said so above, with the reasons. The manuscript is materially better for both.

\end{document}
